\section{Conclusion and Future Work}

This paper presented a decoupled dual pipeline control framework
for quadruped locomotion that achieves robust postural stability
through a novel multistage signal conditioning approach. The
key contributions, namely quintic trajectory planning with
enforced boundary conditions, gain scheduled PID with back
calculation antiwindup, phase aware rotational stance compensation,
and strict architectural isolation at the node level, collectively
enable stable blind locomotion on embedded hardware without
requiring computationally expensive MPC \cite{dicarlo2018dynamic}
or RL solvers.

The proposed multistage pipeline (EMA $\rightarrow$ deadband
$\rightarrow$ gain scheduled PID $\rightarrow$ saturation
$\rightarrow$ LPF) demonstrates that systematic cascading of
conventional signal processing elements can produce corrective
outputs of sufficient quality to stabilize complex mechanical
systems with multiple articulated links, representing a practical
alternative to monolithic optimal controllers.

Future work will pursue three principal directions:
\begin{enumerate}
    \item Incorporation of Ground Reaction Force (GRF) estimation
    via motor current sensing, enabling terrain classification
    without dedicated sensors for force and torque measurement.

    \item Extension of the gait library to incorporate adaptive
    stride modulation, dynamically adjusting $L_{stride}$ and
    $h_{lift}$ based on estimated terrain roughness.

    \item Migration of the gain scheduled PID to an adaptive
    controller that is model free and automatically tunes $K_P$,
    $K_I$, $K_D$ based on performance metrics from the closed
    loop.
\end{enumerate}
